\documentclass[11pt]{article}

\usepackage[margin=1in]{geometry}
\usepackage[T1]{fontenc}
\usepackage[utf8]{inputenc}
\usepackage{setspace}
\usepackage{graphicx}
\usepackage{booktabs}
\usepackage{tabularx}
\usepackage{array}
\usepackage{float}
\usepackage{amsmath}
\usepackage{hyperref}
\usepackage{caption}
\usepackage{enumitem}
\usepackage{xcolor}
\usepackage{longtable}

\setstretch{1.28}
\hypersetup{
    colorlinks=true,
    linkcolor=blue,
    citecolor=blue,
    urlcolor=blue
}

\title{\textbf{LocalLens: Grounded Local-Intelligence Retrieval for Travelers and Recent Movers}}
\author{
\textbf{Instructor: Prof Karl Ni} \\[1em]
Anusha Ravi Kumar \\
\texttt{ravikumar.anu@northeastern.edu}
\and
Sohan Mahesh \\
\texttt{mahesh.so@northeastern.edu}
\and
Parvathi Gottumukkala \\
\texttt{gottumukkala.pa@northeastern.edu}
\and
Shraavya Kishan Bharadwaj \\
\texttt{bharadwaj.sh@northeastern.edu}
}
\date{CS 6120 \quad Natural Language Processing \\ Northeastern University}

\begin{document}
\maketitle

\section{Objectives and Significance}
Travel questions sound simple, but practical local questions are usually much harder than generic web search suggests. A user asking ``Where should I watch the sunset in San Francisco?'' or ``What should I know before moving to Seattle?'' is not only looking for landmark names. They are asking for neighborhood context, tradeoffs, transit practicality, place quality, and the sort of advice that long-time residents learn gradually. General-purpose language models often answer these questions fluently, but not always faithfully. Search engines return many links, but they do not usually synthesize them into a grounded, concise answer. LocalLens was built to address this gap.

The core goal of LocalLens is to create a local-first retrieval augmented generation (RAG) system for travelers and recent movers. The system combines narrative city knowledge, community-style guidance, and structured place records into one answering pipeline. Instead of producing unsupported suggestions from parametric memory alone, the system retrieves evidence from an external corpus, ranks it, and then generates a short readable answer grounded in those retrieved sources.

The significance of the project is twofold. First, it addresses a realistic NLP systems problem: how to answer local questions in a way that is both useful and verifiable. Second, it integrates multiple forms of data that are often treated separately: city guides, background context, community discussion, and structured place metadata. In practice, a useful local assistant needs all of them. A place-only system misses context. A text-only system misses exact place constraints. A model-only system risks hallucination. LocalLens studies how these pieces work together.

\section{Problem Setting}
We frame LocalLens as a grounded local question-answering system over a geographically organized corpus. Given a query $q$, the system must identify:
\begin{itemize}[leftmargin=*]
    \item the relevant location, if one is stated or can be inferred,
    \item the type of intent, such as broad local guidance, place lookup, activity recommendation, newcomer advice, or local custom,
    \item any structured constraints, such as rating thresholds, category filters, or radius/distance conditions,
    \item and the evidence needed to answer the query with minimal hallucination.
\end{itemize}

This means the task is not a single homogeneous retrieval problem. Some questions are primarily narrative:
\begin{quote}
What hidden gems or locals-only spots should I know in San Francisco?
\end{quote}
Other questions are primarily structured:
\begin{quote}
Best taco place in San Jose over 4.5.
\end{quote}
Others are hybrid:
\begin{quote}
Can I rely on public transit in Chicago?
\end{quote}
These require both descriptive text and structured evidence. The system therefore must support multiple retrieval modes while still producing a unified user-facing answer.

\section{Background}
LocalLens is motivated by modern retrieval augmented generation systems for knowledge-intensive tasks. The main idea behind RAG is that a language model should not answer from memory alone when task-specific evidence can be retrieved dynamically. In local-intelligence settings, this is even more important because place quality, neighborhood advice, and visitor guidance depend heavily on external knowledge rather than general pretraining.

This project also draws on ideas from hybrid retrieval. Exact lexical retrieval is useful for neighborhood names, cuisine labels, and transit systems, while dense retrieval is useful when the user and corpus describe the same concept with different wording. LocalLens therefore combines both. It also borrows from practical recommendation systems by keeping a structured place layer in addition to passage retrieval. This allows the system to answer both ``tell me what locals know'' and ``find me a place matching these constraints'' within one application.

A final motivating theme is abstention. In many local-question settings, the best system behavior is not to force an answer. If the corpus does not contain grounded evidence for horse-riding venues near San Jose, the system should say so. This is more trustworthy than returning a semantically similar but incorrect suggestion.

\section{System Overview}
LocalLens is an end-to-end local-intelligence stack rather than a standalone LLM prompt. The project has five connected layers:
\begin{enumerate}[leftmargin=*]
    \item \textbf{Ingestion}: fetch city guides, background summaries, local-discussion signals, curated local-web pages, and structured places.
    \item \textbf{Normalization and storage}: map those heterogeneous sources into a common SQLite-backed schema.
    \item \textbf{Indexing}: chunk long documents and build a dense retrieval index over passage text.
    \item \textbf{Hybrid retrieval and ranking}: combine lexical, dense, and structured place retrieval based on query intent.
    \item \textbf{Grounded answer generation}: produce a readable answer, explanation, and citations through the Streamlit interface.
\end{enumerate}

\begin{figure}[H]
\centering
\fbox{
\parbox{0.94\linewidth}{
\centering
\textbf{Open sources} \\
Wikivoyage, Wikipedia, Reddit/forum digests, curated local-web pages, OpenStreetMap / Overpass, optional Google Places \\[0.5em]
$\downarrow$ \\[0.5em]
\textbf{Normalization layer} \\
source\_documents + chunks + places in SQLite \\[0.5em]
$\downarrow$ \\[0.5em]
\textbf{Dense and lexical indexing} \\
all-MiniLM-L6-v2 embeddings + BM25-style term retrieval \\[0.5em]
$\downarrow$ \\[0.5em]
\textbf{Query planner} \\
structured, semantic, or hybrid route based on location, topic, and constraints \\[0.5em]
$\downarrow$ \\[0.5em]
\textbf{Answer composer} \\
retrieved evidence + local Ollama generation + citations + place cards \\[0.5em]
$\downarrow$ \\[0.5em]
\textbf{Streamlit frontend} \\
answer, why this recommendation, key tips, retrieval notes, citations
}
}
\caption{High-level LocalLens workflow.}
\label{fig:pipeline}
\end{figure}

% FIGURE SLOT 1: SYSTEM ARCHITECTURE / PIPELINE OVERVIEW
% Replace this placeholder with an architecture image if desired.
% Recommended file: docs/figures/source_pipeline_overview.png
\begin{figure}[H]
\centering
\fbox{\parbox{0.88\textwidth}{
\centering
\vspace{1.1cm}
\textbf{Insert Figure 1 Here}\\[0.5em]
System architecture or pipeline overview\\
Suggested file: \texttt{docs/figures/source\_pipeline\_overview.png}
\vspace{1.1cm}
}}
\caption{Optional visual summary of the LocalLens architecture. This is a good place for a cleaned-up version of the slide system-flow graphic.}
\label{fig:pipeline-visual}
\end{figure}

\section{Data Sources and Corpus Construction}
\subsection{Source Types}
The rebuilt LocalLens corpus combines five source types.

\paragraph{Wikivoyage guide documents.}
These are sectioned travel-guide documents for activities, food, lodging, transit, etiquette, and orientation. They are useful because they provide reasonably clean, location-specific prose organized by section.

\paragraph{Wikipedia background documents.}
These provide high-level city and park summaries as orientation context. They are less local-opinion-heavy than guide pages, but useful for anchoring broad descriptions and landmarks.

\paragraph{Forum digests.}
The system ingests local-discussion style documents derived from Reddit threads. This source type aims to capture lived experience and local phrasing that do not appear in formal guides. In practice, these fetches are somewhat brittle because public Reddit endpoints can time out.

\paragraph{Structured place records.}
OpenStreetMap and Overpass contribute restaurants, parks, museums, transit nodes, viewpoints, attractions, hotels, venues, and other place entities. These are kept as structured place rows and also transformed into citation-ready place documents so they can participate in passage retrieval.

\paragraph{Curated local-web / local-guide documents.}
One of the largest late-stage improvements in the project was the addition of a curated local-web ingestion layer. The earlier corpus had strong place coverage but weak long-form local knowledge. The new local-web layer crawls official and curated city pages, scores them for quality, and materializes higher-signal local-guidance documents. When web crawling is sparse, the system now synthesizes fallback \texttt{local\_guide} documents from existing guide, background, and forum sources, which prevents the local-knowledge layer from collapsing to zero.

\subsection{Current Corpus Size}
After the latest full rebuild, the corpus contains the following totals across 24 locations:
\begin{itemize}[leftmargin=*]
    \item \textbf{16,379} source documents
    \item \textbf{21,076} citation-ready chunks
    \item \textbf{14,891} structured place records
\end{itemize}

Table \ref{tab:source-breakdown} shows the source-document breakdown.

\begin{table}[H]
\centering
\caption{Current source-document composition after the latest rebuild.}
\label{tab:source-breakdown}
\begin{tabular}{lr}
\toprule
Source type & Count \\
\midrule
place\_record & 14,891 \\
guide & 944 \\
local\_guide & 208 \\
background & 184 \\
forum\_digest & 152 \\
\bottomrule
\end{tabular}
\end{table}

The large number of place records means the corpus is still structurally place-heavy. However, the presence of 208 \texttt{local\_guide} documents is a major improvement over earlier rebuilds that produced zero such documents. Cities like Seattle, San Jose, Austin, Chicago, and New Orleans now all contain local-guide coverage, which matters for questions about newcomer advice, neighborhood vibe, and local norms.

% FIGURE SLOT 2: SOURCE BREAKDOWN PLOT
% Replace this placeholder with a bar chart of source-document counts.
% Recommended file: docs/figures/source_breakdown.png
\begin{figure}[H]
\centering
\fbox{\parbox{0.8\textwidth}{
\centering
\vspace{0.9cm}
\textbf{Insert Figure 2 Here}\\[0.5em]
Source-document composition plot\\
Suggested file: \texttt{docs/figures/source\_breakdown.png}
\vspace{0.9cm}
}}
\caption{Recommended plot: source-document composition of the rebuilt corpus. This figure would visually reinforce the balance between place records and narrative local-knowledge sources.}
\label{fig:source-breakdown-plot}
\end{figure}

\subsection{Place Coverage}
Table \ref{tab:place-cats} shows the largest structured categories in the current corpus.

\begin{table}[H]
\centering
\caption{Largest place categories in the rebuilt corpus.}
\label{tab:place-cats}
\begin{tabular}{lr}
\toprule
Category & Count \\
\midrule
restaurant & 3,920 \\
park & 2,304 \\
attraction & 2,160 \\
hotel & 1,770 \\
transit & 1,761 \\
venue & 1,334 \\
museum & 1,172 \\
viewpoint & 263 \\
market & 138 \\
\bottomrule
\end{tabular}
\end{table}

This table also explains several later failure modes. LocalLens has abundant restaurant, park, attraction, and transit-stop coverage, but much thinner coverage for rare categories like equestrian or dark-sky activity-specific venues. As a result, some queries are answered well, while others appropriately abstain.

% FIGURE SLOT 3: PLACE CATEGORY DISTRIBUTION
% Replace this placeholder with a plot of top place categories.
% Recommended file: docs/figures/place_category_distribution.png
\begin{figure}[H]
\centering
\fbox{\parbox{0.84\textwidth}{
\centering
\vspace{0.9cm}
\textbf{Insert Figure 3 Here}\\[0.5em]
Place-category distribution plot\\
Suggested file: \texttt{docs/figures/place\_category\_distribution.png}
\vspace{0.9cm}
}}
\caption{Recommended plot: largest structured place categories in the corpus. This figure helps explain why LocalLens is strong on some query families and weaker on others.}
\label{fig:place-category-plot}
\end{figure}

\section{Storage, Indexing, and Retrieval}
\subsection{SQLite Schema}
The project uses SQLite as its primary local data store. The core tables are:
\begin{itemize}[leftmargin=*]
    \item \texttt{source\_documents}: long-form city documents, forum digests, local-guide syntheses, and place-derived mini-documents.
    \item \texttt{chunks}: passage-level segments used for citation and dense retrieval.
    \item \texttt{places}: structured place records used for exact filtering and place-card ranking.
\end{itemize}

This design is simple to rebuild, easy to ship, and consistent with the reproducibility goals of the course project.

\subsection{Dense Indexing}
Each chunk is embedded with \texttt{sentence-transformers/all-MiniLM-L6-v2}. The implementation supports a Qdrant-backed vector path, but it also falls back gracefully to a local NumPy embedding matrix when Qdrant is unavailable. In the latest successful rebuild, the dense index was written as a local artifact:
\begin{quote}
\texttt{artifacts/chunk\_embeddings.npy}
\end{quote}
using the backend label \texttt{numpy:sentence-transformers/all-MiniLM-L6-v2}. This fallback kept the project stable during rebuild and deployment.

\subsection{Hybrid Retrieval}
Retrieval combines three signals:
\begin{itemize}[leftmargin=*]
    \item lexical retrieval over chunk text,
    \item dense retrieval over embeddings,
    \item structured place search over exact fields.
\end{itemize}

For passage fusion, LocalLens uses a reciprocal-rank style merging intuition:
\[
\mathrm{score}(d) = \sum_{r \in R} \frac{1}{k + \mathrm{rank}_r(d)}
\]
where $R$ denotes the participating ranked lists and $k$ is a smoothing constant. In practice, this improves coverage because dense retrieval recovers semantic matches while lexical retrieval preserves exact city/topic terms.

\subsection{Query Planning}
A major part of the project is the query planner implemented in the service layer. Each user query is parsed for:
\begin{itemize}[leftmargin=*]
    \item location,
    \item topic,
    \item activity type,
    \item rating constraints,
    \item audience hints,
    \item vibe or local-knowledge intent,
    \item and distance/radius language such as ``under 2 hours from San Jose.''
\end{itemize}

The planner then routes the query as one of three modes:
\begin{itemize}[leftmargin=*]
    \item \textbf{structured}: primarily place filtering and ranking,
    \item \textbf{semantic}: primarily local-knowledge passage retrieval,
    \item \textbf{hybrid}: combine both.
\end{itemize}

This hybrid planner is one of the clearest ways the system evolved. Early versions were too eager to answer all questions through loosely matched places. The newer system is better at separating questions like ``best taco place'' from ``what should a newcomer know before moving to Seattle?'' even if the final ranking is still imperfect.

\section{Curated Local-Web Ingestion Improvements}
One of the most important late-stage modifications was the introduction of a stronger local-web pipeline. Earlier versions of the project relied too heavily on place records and guide documents. That made the system broad, but not sufficiently local in the ``people who live there would know this'' sense.

The new local-web layer works as follows:
\begin{enumerate}[leftmargin=*]
    \item Start from official and curated city URLs.
    \item Discover additional same-domain links using anchor text and URL-path keywords.
    \item Expand discovery through sitemap parsing and limited second-hop exploration.
    \item Score pages for quality, topic fit, and local specificity rather than simply keeping the first fixed number of pages.
    \item Select the strongest local pages per city/topic/domain.
    \item If that still yields thin results, synthesize \texttt{local\_guide} documents from existing guide, background, and forum sources.
\end{enumerate}

This matters because it shifts the project away from a ``small hard cap of arbitrary pages'' to a quality-scored local-knowledge layer. It also aligns better with the user-facing goal: helping someone understand a place more like a local rather than only surfacing POIs.

\section{Generation and Abstention}
LocalLens is designed to use a local Ollama-served model for grounded answer generation. The prompt layer asks the model to:
\begin{itemize}[leftmargin=*]
    \item stay grounded in retrieved evidence,
    \item produce a readable narrative answer,
    \item explain why the recommendation was made,
    \item provide concise practical tips,
    \item and abstain when evidence is weak.
\end{itemize}

The abstention behavior turned out to be especially important. A good example is:
\begin{quote}
\texttt{horse riding activities under 2 hours from san jose}
\end{quote}
Earlier versions of the system sometimes answered such a query by latching onto unrelated place names containing words like ``stable.'' The current version correctly says that it cannot find grounded equestrian evidence in the corpus. This is not a glamorous answer, but it is substantially more trustworthy.

\section{Frontend and User Experience}
The frontend is a Streamlit application that presents several layers of evidence instead of a single opaque answer. The main interface includes:
\begin{itemize}[leftmargin=*]
    \item a free-form query box,
    \item optional location/topic filters,
    \item a narrative answer section,
    \item a \emph{Why This Recommendation} section,
    \item key tips,
    \item retrieval notes,
    \item place cards,
    \item and clickable citations.
\end{itemize}

The current public deployment used for demo testing is hosted at:
\begin{quote}
\url{http://136.117.92.234:8501}
\end{quote}

This interface matters because LocalLens is not only about retrieving facts. It is about making those facts useful and readable to a traveler or recent mover.

% FIGURE SLOT 4: MAIN APP SCREENSHOT
% Replace this placeholder with a real screenshot from the LocalLens app.
% Recommended file: docs/figures/app_main_ui.png
\begin{figure}[H]
\centering
\fbox{\parbox{0.92\textwidth}{
\centering
\vspace{1.1cm}
\textbf{Insert Screenshot 1 Here}\\[0.5em]
Main Streamlit interface with query box, answer, retrieval notes, and citations\\
Suggested file: \texttt{docs/figures/app\_main\_ui.png}
\vspace{1.1cm}
}}
\caption{Recommended screenshot: the main LocalLens interface with a representative answer and retrieval notes visible.}
\label{fig:ui-main}
\end{figure}

\section{Rubric Alignment and Deliverables}
The project was also shaped around the practical expectations of the CS 6120 final-project rubric rather than only around a research-style architecture diagram. Table \ref{tab:rubric} maps the current implementation to the main deliverables emphasized in the repository and demo materials.

\begin{table}[H]
\centering
\caption{Rubric-alignment summary for the current LocalLens submission.}
\label{tab:rubric}
\small
\begin{tabularx}{\textwidth}{>{\raggedright\arraybackslash}p{0.28\textwidth} >{\raggedright\arraybackslash}X}
\toprule
Rubric-aligned requirement & How LocalLens satisfies it \\
\midrule
Frontend / usable interface & Streamlit frontend with a free-form prompt box, answer panel, why-this-recommendation explanation, key tips, retrieval notes, place cards, and clickable evidence blocks. \\
10k+ entry data layer & Latest rebuild contains 16,379 source documents, 21,076 chunks, and 14,891 structured place records across 24 locations. \\
Local model usage & Generation path is designed around a locally served Ollama model, and the public deployment points the app to a local Ollama endpoint rather than a hosted proprietary API. \\
Grounded citations & Every answer can surface passage-level citations with source URL, location, topic, and retrieved text so the user can inspect the evidence directly. \\
Reproducibility & The pipeline is rebuilt with \texttt{build\_corpus.py}, \texttt{build\_index.py}, and \texttt{run\_presentation\_eval.py}, with SQLite and artifact files stored locally. \\
Deployment / engineering completeness & The repository includes a Dockerfile, docker-compose configuration, README instructions, GCP notes, and a working public Streamlit endpoint for demo testing. \\
Analysis and documentation & The writeup documents architecture, corpus composition, retrieval design, evaluation findings, failure modes, and future work, rather than only presenting a product overview. \\
\bottomrule
\end{tabularx}
\end{table}

This alignment matters because the project is not only intended to be technically interesting. It is also meant to be inspectable, runnable, and demoable under course constraints. In particular, LocalLens satisfies the important systems-oriented expectations of the rubric: a frontend, a sufficiently large database, local model support, explicit citations, and a reproducible build path.

\section{Evaluation Setup}
To evaluate LocalLens, we created a small presentation-oriented benchmark of ten realistic local questions spanning sunsets, hidden gems, newcomer advice, neighborhoods, food, transit, local norms, nightlife, niche activities, and family-friendly outdoor activities. These queries are executed through a script that prints the answer text, explanation, source summary, citation types, and place-card count.

The benchmark is not meant to be a statistically complete evaluation dataset. Instead, it is a practical stress test that reflects the kinds of questions instructors or demo viewers actually ask. We used it as a presentation-safety benchmark rather than as a purely offline ML benchmark, because the classroom risk is not only average retrieval quality but whether the system behaves sensibly on realistic public-facing questions.

We therefore evaluate along four axes:
\begin{itemize}[leftmargin=*]
    \item \textbf{Groundedness}: whether the retrieved citations actually support the answer text.
    \item \textbf{Local usefulness}: whether the answer surfaces practical guidance instead of generic city descriptions.
    \item \textbf{Structured precision}: whether exact constraints such as rating, category, and location are respected.
    \item \textbf{Failure behavior}: whether the system abstains honestly when evidence is missing instead of fabricating a plausible but unsupported answer.
\end{itemize}

This evaluation framing is important because LocalLens is ultimately a product-like NLP system. A technically sophisticated retriever is not sufficient if the user-facing answer is misleading, and a polished interface is not sufficient if the answer is ungrounded.

\section{Results and Analysis}
\subsection{What Improved}
The latest rebuild showed several clear improvements relative to earlier versions.

\paragraph{Local-guide coverage now exists.}
The new corpus contains 208 \texttt{local\_guide} documents. This is a meaningful change, because earlier rebuilds produced zero such documents and therefore left newcomer-advice and hidden-gem questions under-supported.

\paragraph{Hidden-gem style questions improved.}
For the San Francisco hidden-gems query, the system now retrieves both \texttt{forum\_digest} and \texttt{local\_guide} sources instead of collapsing to generic tourism or unrelated place-card content.

\paragraph{Safer abstention works.}
The horse-riding query no longer hallucinates. It explicitly reports that there is no grounded equestrian evidence for that request in the current corpus.

\paragraph{Scenic and family/outdoor questions are better grounded.}
Queries like ``Good sunset spots in San Francisco'' and ``What are some family-friendly outdoor things to do in San Diego?'' now pull from a mix of \texttt{local\_guide}, \texttt{background}, and \texttt{place\_record} evidence.

These improvements matter because they reflect a real change in system behavior, not just a larger corpus size. Earlier project iterations often looked broad on paper because they had thousands of place rows, but broad place coverage alone did not reliably answer the kinds of questions a traveler or newcomer actually asks. The current rebuild demonstrates that adding a stronger local-guide layer changes the retrieval mix itself. Instead of answering most queries through generic POI matches, the system now has a better chance of grounding answers in local-knowledge documents, especially for ``hidden gems,'' broad sightseeing, and some activity-style questions.

\subsection{Interpreting the Rebuilt Corpus}
The current corpus composition helps explain both the system's strengths and its remaining weaknesses. On the one hand, 208 \texttt{local\_guide} documents are enough to noticeably improve certain long-form local questions. On the other hand, the corpus still contains 14,891 place-derived documents. This means that the dominant evidence type in the system is still place-centric rather than discussion-centric. In practice, this creates a tension between breadth and relevance. The system knows that many places exist, but it is still learning when not to let place presence dominate the final answer.

This also explains why the project improved first on scenic and broad discovery questions before it improved on local-norm and neighborhood-vibe questions. A scenic question such as ``good sunset spots'' maps naturally to viewpoints, parks, waterfronts, and attractions that already exist in the corpus. A question such as ``what should a newcomer know before moving to Seattle'' depends on a much subtler evidence layer: local customs, mobility tradeoffs, neighborhood tone, and practical adaptation advice. Those patterns are not captured well by place records alone and require richer local prose.

\subsection{Representative Evaluation Outcomes}
Table \ref{tab:eval} summarizes the latest representative results.

\begin{table}[H]
\centering
\caption{Representative evaluation outcomes from the latest rebuilt corpus.}
\label{tab:eval}
\small
\begin{tabularx}{\textwidth}{>{\raggedright\arraybackslash}p{0.28\textwidth} >{\raggedright\arraybackslash}p{0.22\textwidth} >{\raggedright\arraybackslash}X}
\toprule
Query & Outcome & Interpretation \\
\midrule
Good sunset spots in San Francisco & Stronger grounded answer & Retrieved from \texttt{background}, \texttt{local\_guide}, and place records; useful improvement. \\
What hidden gems or locals-only spots should I know in San Francisco? & Improved local-knowledge retrieval & Uses \texttt{forum\_digest} and \texttt{local\_guide}; much closer to the intended product behavior. \\
What should a newcomer know before moving to Seattle? & No grounded evidence returned & Local-guide coverage exists, but retrieval/routing is still not consistently surfacing the right evidence. \\
Which neighborhoods in Seattle feel quiet but still interesting? & No grounded evidence returned & Indicates neighborhood-vibe retrieval remains weak. \\
Best taco place in San Jose over 4.5 & Incorrect ranking & Returned Taco Bell, revealing that exact filtering and ranking need further tightening. \\
Can I rely on public transit in Chicago? & Weak structured overreach & Returned stop-level place matches rather than a stronger transit narrative. \\
What local norms should I know before visiting New Orleans? & No grounded evidence returned & Local-custom coverage remains sparse or under-ranked. \\
Where do locals actually go for nightlife in Austin? & Noisy venue ranking & Nightlife retrieval still spills into generic venue/place matches. \\
horse riding activities under 2 hours from san jose & Correct abstention & Properly reports lack of grounded equestrian evidence. \\
What are some family-friendly outdoor things to do in San Diego? & Moderately improved & Grounded in beach- and outdoor-related evidence, though still somewhat place-heavy. \\
\bottomrule
\end{tabularx}
\end{table}

An important detail from this evaluation is that the local evaluation run reported \texttt{used\_local\_llm = False} for all ten benchmark queries. That means these results expose the retrieval and ranking behavior more than the final live-generation fluency. This is helpful diagnostically: it tells us the current main bottleneck is not just the generator, but the retrieval and selection logic underneath it.

\subsection{Detailed Query-Family Analysis}
The benchmark is most informative when grouped by query family instead of viewed as ten isolated prompts.

\paragraph{1. Place-discovery questions.}
The system is currently strongest when the query can be anchored to a place category that is already well represented in the corpus. The San Francisco sunset query and the San Diego family-outdoors query are both examples of this. In these cases, LocalLens benefits from the combination of structured category coverage, scenic-place metadata, and supplementary local-guide text. The answer is not perfect, but it is directionally useful and grounded.

At the same time, the San Jose taco result shows that place-discovery quality is not solved simply by having many restaurants in the database. The system correctly interprets the query as restaurant-focused and rating-sensitive, but the ranking stack still allows low-quality or chain-like candidates to surface too highly. This is a classic retrieval-systems lesson: recall is not the same as precision. The corpus contains enough restaurant data, but the decision layer above it still needs stronger filtering and better ranking features.

\paragraph{2. Local-knowledge narrative questions.}
The hidden-gems query in San Francisco is one of the clearest signs that the project moved in the right direction. Earlier builds often responded to this style of prompt with generic tourism passages. In the latest evaluation, the system retrieves from \texttt{forum\_digest} and \texttt{local\_guide}, which means the new ingestion path is actually influencing the answer source mix. This is exactly the intended product behavior: if a user asks for ``locals-only'' knowledge, the evidence should come disproportionately from discussion and local-guide style sources rather than from encyclopedic city summaries.

However, Seattle newcomer advice and Seattle neighborhood-vibe questions are still weak. This suggests that the system has not yet reached the point where local-guide coverage alone is sufficient. The likely issue is that these topics require both source depth and better routing. The documents exist in at least some form, but the retrieval stack is not yet reliably surfacing them for those prompts.

\paragraph{3. Norms, etiquette, and newcomer-friction questions.}
The New Orleans local-norms query failing with no grounded evidence is revealing. These questions are among the hardest in the entire project because they depend on advice that is often informal, city-specific, and unevenly documented. Unlike restaurants or parks, there is no stable structured schema for ``local norms.'' This means LocalLens must depend on high-quality prose sources and semantic matching. The current system architecture is capable of that in principle, but the corpus and ranking layer are not yet strong enough in practice.

\paragraph{4. Transit and logistics questions.}
The Chicago transit result illustrates a subtle but important failure mode. The system clearly recognizes transit as the relevant topic, yet still returns stop-level place matches rather than a city-level answer about how usable the transit system actually is. This means the issue is not missing transit data but evidence prioritization. For this query family, the system should prioritize textual guidance about reliability, coverage, and tradeoffs over bare stop entities. In other words, the error is architectural but fixable: the right evidence is not necessarily absent, but it is being outcompeted by overly literal structured matches.

\paragraph{5. Niche-activity questions.}
The horse-riding query is a strong example of a system failing correctly. The current corpus has only minimal equestrian coverage, so the honest answer is abstention. This result is useful because it shows that the system's safety behavior improved. In a live demo or class setting, this is better than confidently inventing a recommendation. It also highlights a key product lesson from the project: a trustworthy local assistant must know when the corpus is insufficient.

% FIGURE SLOT 5: QUALITATIVE QUERY EXAMPLES
% Replace this placeholder with a montage of strong answer + abstention examples.
% Recommended file: docs/figures/eval_query_examples.png
\begin{figure}[H]
\centering
\fbox{\parbox{0.92\textwidth}{
\centering
\vspace{1.1cm}
\textbf{Insert Figure 5 Here}\\[0.5em]
Qualitative query output examples\\
Suggested file: \texttt{docs/figures/eval\_query\_examples.png}
\vspace{1.1cm}
}}
\caption{Recommended qualitative figure: representative query outputs from the live app, showing both a successful grounded answer and a correctly abstained answer.}
\label{fig:eval-examples}
\end{figure}

\subsection{Failure Analysis}
The evaluation highlights several persistent failure modes that are important both technically and from a grading perspective, because they show that we analyzed the system as an end-to-end NLP pipeline rather than only celebrating the successful examples.

\paragraph{Place-heavy bias.}
The corpus is still dominated by place-derived documents. This can be useful, but it also means the system is tempted to answer many questions through place cards even when a narrative local-guidance answer would be more appropriate.

\paragraph{Ranking is still too eager in some structured queries.}
The San Jose taco query returning Taco Bell is the clearest example. The system recognizes the request as a food/place query, but the ranking logic still needs stronger chain downweighting, stricter rating use, and better cuisine matching.

\paragraph{Some local-knowledge topics remain thin.}
Questions about Seattle newcomer advice, Seattle neighborhood vibe, and New Orleans local norms are still not reliably grounded. This suggests that simply adding \texttt{local\_guide} documents is not enough; the retrieval and topic mapping for these questions must also improve.

\paragraph{Source fragility propagates into uneven coverage.}
The logs from the latest rebuild show that Reddit/forum fetches can fail by city because of public endpoint timeouts. The corpus build continues, which is good for robustness, but the resulting data quality becomes uneven. A city can therefore have strong guide and place coverage but weak local-discussion evidence, which directly affects the kinds of answers LocalLens can produce there.

\paragraph{Dense retrieval infrastructure is more flexible than fully standardized.}
The project was designed with a vector-database path in mind, but the latest stable rebuild relied on the NumPy artifact fallback rather than a consistently active Qdrant backend. Functionally this still works, but analytically it means the deployment is currently optimized more for stability and reproducibility than for final-scale serving architecture.

\paragraph{Evaluation currently emphasizes breadth and realism more than statistical benchmarking.}
The present benchmark is intentionally realistic, but it is still small. It is very useful for presentation safety and system debugging, yet it does not replace a larger annotated evaluation set with query labels, groundedness judgments, and answer-quality scoring. This is not a flaw in the system itself, but it is a limitation in how comprehensively we can quantify progress right now.

% FIGURE SLOT 6: QUERY FAMILY SUMMARY
% Replace this placeholder with a summary plot of strong / partial / weak query families.
% Recommended file: docs/figures/query_family_summary.png
\begin{figure}[H]
\centering
\fbox{\parbox{0.86\textwidth}{
\centering
\vspace{0.9cm}
\textbf{Insert Figure 6 Here}\\[0.5em]
Query-family summary plot\\
Suggested file: \texttt{docs/figures/query\_family\_summary.png}
\vspace{0.9cm}
}}
\caption{Recommended plot: qualitative summary of current performance by query family, such as scenic, hidden-gem, newcomer, transit, nightlife, and niche-activity prompts.}
\label{fig:query-family-summary}
\end{figure}

\section{Deployment and Reproducibility}
LocalLens is designed to be reproducible end to end. The main rebuild path is:
\begin{verbatim}
python3 scripts/build_corpus.py
python3 scripts/build_index.py
python3 scripts/run_presentation_eval.py
PYTHONPATH=src streamlit run app.py
\end{verbatim}

For public demo testing, the project was also pushed to a remote VM and the Streamlit application was restarted there after the latest rebuild and index creation. The deployment uses Docker-compatible project structure, local artifact files for the dense index, and a local Ollama endpoint for generation.

The repository is available at:
\begin{quote}
\url{https://github.com/anusha24-bit/LocalLens}
\end{quote}

\section{Limitations}
LocalLens still has several limitations, and many of them are structural rather than cosmetic. These limitations matter because they explain why the system can already be compelling in a live demo while still falling short of being a generally reliable local-intelligence engine.

\begin{itemize}[leftmargin=*]
    \item \textbf{City coverage quality is uneven even when corpus size is large.} A database with more than 16,000 documents can still be weak on the exact kind of advice users want. This is especially true for neighborhood-vibe, newcomer-friction, and local-custom questions, which require richer narrative evidence than formal travel-guide prose usually provides.
    \item \textbf{Reddit/forum ingestion is brittle.} Public Reddit endpoints can timeout or return empty results. Because the build intentionally continues even when this happens, some cities end up with uneven forum-style coverage. This directly affects the quality of hidden-gem, local-opinion, and local-norm answers.
    \item \textbf{The corpus is still structurally place-heavy.} The majority of stored documents are place-derived. This helps with broad coverage, but it also biases ranking toward POIs, addresses, and category matches even in cases where the better answer would come from local narrative evidence.
    \item \textbf{Structured ranking is not yet sophisticated enough.} The system can identify that a query is restaurant-, transit-, or nightlife-related, but it does not yet use a mature enough feature set to consistently choose the best answer among many candidate entities. Cuisine specificity, chain downweighting, and contextual desirability are still weak.
    \item \textbf{Regional distance queries are only partially solved.} Queries like ``under 2 hours from San Jose'' require more than text similarity. They depend on geospatial reasoning, approximate drive-time logic, and regional candidate expansion, all of which are only partially handled in the current version.
    \item \textbf{The active dense retrieval backend prioritized robustness over architectural ambition.} Although the code supports a Qdrant-backed vector path, the latest stable deployment used the NumPy-based dense artifact. This was the correct engineering tradeoff for stability, but it means the system is not yet demonstrating its final intended retrieval infrastructure.
    \item \textbf{Evaluation depth is still limited.} The current evaluation is highly useful for product-level inspection, but it is not yet a large labeled benchmark. We therefore know a lot about visible failure cases, but less about overall distributional performance across many hundreds of queries.
\end{itemize}

Taken together, these limitations show that LocalLens is already a credible systems project but not yet a fully mature local-advice engine. That distinction is important for an honest report: the project succeeds as a thoughtful NLP system, but it also reveals how difficult local groundedness really is.

\section{Future Work}
There are several high-value next steps, and they fall into three broad groups: stronger data, stronger ranking, and stronger evaluation.

\begin{itemize}[leftmargin=*]
    \item \textbf{Build richer local-knowledge sources.} The single biggest future gain would come from increasing high-quality narrative sources for newcomer advice, neighborhood tone, local etiquette, and ``what locals actually do'' style questions. The current local-web layer is a strong start, but it should be expanded with better curated source registries and more robust extraction.
    \item \textbf{Replace fragile discussion ingestion with a more reliable pipeline.} Forum-style local advice is clearly valuable, but relying on brittle public Reddit endpoints makes the corpus inconsistent. A future version should move to a more dependable, policy-compliant ingestion path or expand into better curated community and city-publication sources.
    \item \textbf{Improve place ranking with domain-aware signals.} Restaurant, nightlife, and transit queries would benefit from stronger rank features such as explicit chain penalties, cuisine specificity, neighborhood desirability, review quality, and mode-specific rules. In other words, the structured layer should not only filter candidates but also reason about what makes a candidate good for the user's intent.
    \item \textbf{Introduce query-family-specific reranking.} The current system still lets certain query families collapse into generic entity matching. Transit questions should privilege system-level transit evidence, nightlife questions should privilege local-opinion and venue-style nightlife signals, and newcomer questions should prefer broad orientation and adaptation advice over attractions.
    \item \textbf{Strengthen geospatial reasoning.} Regional queries such as ``under two hours from San Jose'' require dedicated support for radius expansion, route approximation, and city-adjacent place discovery. A future version should incorporate more explicit geographic search logic rather than depending mostly on textual matching.
    \item \textbf{Standardize the vector backend.} The retrieval layer currently values stability and reproducibility, which was the right choice for submission, but a future version should settle on one consistently active dense-index path and benchmark it carefully. That includes comparing NumPy artifacts and Qdrant-backed serving under the same corpus and query suite.
    \item \textbf{Create a larger annotated evaluation set.} A stronger research-grade follow-up would include a labeled benchmark of local questions with annotations for groundedness, local usefulness, citation support, and abstention correctness. That would make progress measurable beyond anecdotal examples.
    \item \textbf{Measure the effect of the final generation model separately from retrieval.} Because some of the local evaluation runs exposed retrieval behavior more than generation behavior, a future experimental pass should separately quantify how much answer quality changes when the exact same retrieved evidence is passed through different local LLMs.
\end{itemize}

These future directions are intentionally concrete. They are not generic ``make it better'' statements; they follow directly from the failure cases surfaced by the rebuilt corpus and the presentation benchmark.

\section{Conclusion}
LocalLens demonstrates that local-question answering is best treated as a hybrid NLP systems problem rather than a pure language-model prompt. Useful local answers require multiple kinds of evidence: guide prose, local-discussion style signals, and structured place metadata. The latest version of the project substantially improved the local-knowledge layer through curated local-web ingestion, quality-based page selection, and fallback local-guide synthesis. It also improved system honesty through better abstention.

At the same time, the project remains transparent about its weaknesses. Some of the most challenging queries still fail because of ranking and source-coverage gaps, not because the interface or model is missing. This honesty is part of the contribution. LocalLens shows both what a grounded local-intelligence system can already do and where careful retrieval engineering still matters.

\section{Contributions}
\begin{itemize}[leftmargin=*]
    \item \textbf{Anusha Ravi Kumar}: contributed across product design, system integration, evaluation, and communication. This included shaping the user-facing direction of LocalLens as a local-intelligence assistant, participating in integration and deployment work, helping drive iterative testing of the live system, and contributing to the final presentation and report materials.
    \item \textbf{Sohan Mahesh}: contributed across retrieval engineering, data iteration, evaluation, and system refinement. This included work on ranking and routing behavior, corpus-build iteration, debugging failure cases surfaced by evaluation, and helping turn those findings into system-level improvements and demo-ready behavior.
    \item \textbf{Parvathi Gottumukkala}: contributed across data engineering, source analysis, evaluation, and communication. This included reviewing source quality, supporting corpus design decisions, helping assess whether retrieved evidence captured real local knowledge, and contributing to the presentation and report-oriented framing of the project.
    \item \textbf{Shraavya Kishan Bharadwaj}: contributed across documentation, evaluation, presentation preparation, and overall project packaging. This included helping translate technical implementation into clearer user-facing and submission-facing artifacts, supporting evaluation and demo preparation, and contributing to the final communication materials around the project.
\end{itemize}

Overall, the project should be understood as an equal-contribution team effort. All four team members contributed to the technical build, iterative evaluation, presentation preparation, and final submission materials, even though particular tasks naturally had moments of primary ownership.

\end{document}
